\section{实施检查与持续改进}
\subsection{常规检查表}
Graham 建议在每次 Agent 迭代前后执行检查表，确保 prompt、工具链、日志和审批机制都处于可用状态。检查表项包括：是否更新 open-source requirements、是否验证镜像权限、是否更新测试基线、是否同步 Slack 通知。

\begin{table}[H]
\centering
\begin{tabular}{p{0.32\textwidth} p{0.6\textwidth}}
\toprule
检查项 & 描述 \\
\midrule
Prompt 状态 & 确保 prompt 模板包含最新的 issue 样本与 guardrails \\
\midrule
工具链更新 & 检查 docker 镜像、依赖版本与 API 权限是否过期 \\
\midrule
日志与审计 & 验证事件流是否正常，并且 audit log 可回放 \\
\midrule
审批机制 & 复核高风险命令审批流程是否完整，包括 emergency override \\
\bottomrule
\end{tabular}
\caption{实施前后的常规检查项。}
\end{table}

\subsection{持续改进循环}
每个 Agent 迭代都应该包含一个回顾环节：回顾成功的案例、失败的经验，并把新的 insight 写入 FAQ/文档，供下一轮模型训练使用。Graham 强调“1% 改进常态化”——每天把指令、提示和工具链中的某一项提一点。

\begin{importantbox}{小步快跑}
把 Agent 更新分解成 5 分钟的改动，验证通过后再合并，尽量避免一次性大量变更。
\end{importantbox}

\begin{knowledgebox}{改进指标}
使用“成功案例比例”、“复盘完成率”和“文档更新次数”作为改进指标，把工程和研究的工作闭环在一起。
\end{knowledgebox}

\subsection{本章小结}
\begin{itemize}
    \item 检查表保障 Agent 架构在每次迭代前后都处于健康状态。
    \item 小步快跑与改进指标让团队可以在风险可控的前提下持续进步。
\end{itemize}

